% =============================================================================
% APPENDIX NNN: METHOD-MULTI-LEVEL
% Multi-Level Implementation Templates and Parameter Tables
% =============================================================================
%
% Code: NNN
% Category: METHOD-MULTI-LEVEL (Multi-level implementation toolkit)
% Language: English
% Version: 1.0 (January 2026)
%
% Purpose: Provide copy-paste-ready diagnostic templates, 9C parameter tables,
% and portfolio design worksheets for each of the five organizational levels:
% Individual (L0), Household (L1), Organization (L2), National (L3), Global (L4).
%
% For: Practitioners implementing programs/policies at any level who need to
% understand how the framework scales and what changes at each level.
%
% Dependencies: Chapter 23 (Multi-Level Implementation), Chapter 16-22 (foundation)
%
% Cross-references: Appendix HHH, KKK, LLL, MMM
%
% =============================================================================

\newpage
\section*{Appendix NNN: METHOD-MULTI-LEVEL}
\label{app:method-multi-level}
\addcontentsline{toc}{section}{Appendix NNN: METHOD-MULTI-LEVEL}

% =============================================================================
% HEADER
% =============================================================================

\begin{tcolorbox}[colback=blue!5!white, colframe=blue!75!black, title=Quick Reference}
\small
\textbf{Purpose:} Understand how 9C framework scales across five organizational levels. Templates provided for each level.

\textbf{Key Principle:} Framework structure stays the same; parameters and constraints vary by level.

\textbf{Five Levels:}
\begin{enumerate}[nosep]
\textbf{Appendix:} NNN (METHOD)

\item \textbf{Level 0: Individual} --- Single person decision-making
\item \textbf{Level 1: Household} --- Couple/family with coupled decisions
\item \textbf{Level 2: Organization} --- Company/department with hierarchy
\item \textbf{Level 3: National} --- Regional/national government policy
\item \textbf{Level 4: Global} --- International coordination without enforcement
\end{enumerate}

\textbf{Process:} Choose your level → Fill diagnostic template → Use parameter table → Adapt portfolio design
\end{tcolorbox}



\begin{abstract}
This appendix provides systematic analysis and documentation relevant to the
Evidence-Based Framework (EBF). It integrates with the 10C CORE architecture
and provides validated findings for practical application.
\end{abstract}

\begin{tcolorbox}[
    colback=orange!5!white,
    colframe=orange!75!black,
    title={\textbf{Appendix Scope: Ziel / In-Scope / Out-of-Scope / Constraints}},
    fonttitle=\bfseries
]

\textbf{Ziel (Objective):}
\begin{itemize}[nosep]
    \item Provide systematic analysis for METHOD integration
\end{itemize}

\textbf{In-Scope:}
\begin{itemize}[nosep]
    \item Core analysis and findings
    \item Framework integration points
\end{itemize}

\textbf{Out-of-Scope:}
\begin{itemize}[nosep]
    \item Detailed empirical replication $\rightarrow$ METHOD appendices
\end{itemize}

\textbf{Constraints:}
\begin{itemize}[nosep]
    \item Analysis bounded by available data
\end{itemize}

\end{tcolorbox}

---

% =============================================================================
% LEVEL 0: INDIVIDUAL
% =============================================================================

\subsection*{Level 0: Individual --- Diagnostic Template}

\textbf{Unit:} Single person making behavioral decision

\textbf{Characteristic:} Variation is \emph{between} individuals (heterogeneous preferences)

\begin{center}
\begin{tabular}{@{}lccl@{}}
\toprule
\textbf{9C Dimension} & \textbf{Your Value} & \textbf{Measure} & \textbf{Data Source} \\
\midrule
L (Who) & [Age, demographics] & Years / Category & Survey \\
d (What) & [Utilities: FEPSDE] & Importance scores & Survey \\
γ (How) & [Key interactions] & Correlation matrix & Survey / Observed \\
Ψ (Context) & [8 dimensions] & Scale 0-1 & Assessment \\
Θ (Parameters) & [Intervention effects] & Percentage improvement & Literature / Pilot \\
A (Awareness) & \_\_\_\% & Subjective probability & Direct question \\
W (Readiness) & \_\_\_\% & Willingness 0-1 & Behavior + intent \\
φ (Phase) & \_\_\_\_ & Stage 0-1 & Questionnaire \\
Scope & Individual & Level of analysis & \\
\midrule
\end{tabular}
\end{center}

\textbf{Example completed:}
\begin{center}
\begin{tabular}{@{}lll@{}}
\toprule
\textbf{Dimension} & \textbf{Value} & \textbf{Interpretation} \\
\midrule
L & Age 52, M, employed & Mid-career \\
d & Highest: Financial (0.8), Health (0.6) & Money and health matter most \\
A & 25\% & Low awareness of climate risk \\
W & 0.3 & Low readiness to change \\
φ & 0.2 (PreAware) & Early stage \\
Segment & Present-Biased (0.7) & Immediate gratification focused \\
\midrule
\end{tabular}
\end{center}

Portfolio design at L0: Use $I_{\text{AWARE}}$ (information on financial benefits of behavior change) + $I_{\text{AWARE},\kappa}$ (regular feedback)

---

% =============================================================================
% LEVEL 1: HOUSEHOLD
% =============================================================================

\subsection*{Level 1: Household --- Diagnostic Template}

\textbf{Unit:} Couple or multi-member household with shared decisions

\textbf{Characteristic:} \emph{Coupled} preferences; asymmetric information; phase mismatch

\subsubsection*{1.1 Household Joint Diagnostic}

\begin{center}
\small
\begin{tabular}{@{}p{2cm}p{2cm}p{2cm}p{2cm}@{}}
\toprule
\textbf{Dimension} & \textbf{Partner A} & \textbf{Partner B} & \textbf{Mismatch?} \\
\midrule
Phase (φ) & \_\_\_\_ & \_\_\_\_ & YES / NO \\
Awareness (A) & \_\_\_\% & \_\_\_\% & Difference \_\_\_pp \\
Readiness (W) & \_\_\_\_ & \_\_\_\_ & Conflicting? \\
Key utility (d) & \_\_\_\_ & \_\_\_\_ & Aligned? \\
\midrule
\end{tabular}
\end{center}

\subsubsection*{1.2 Household Complementarity Matrix}

\begin{quote}
Estimate γ (complementarity) for key decision pairs:
\end{quote}

\begin{center}
\begin{tabular}{@{}llc@{}}
\toprule
\textbf{Decision Pair} & \textbf{Description} & \textbf{γ Estimate} \\
\midrule
A's commitment + B's buy-in & One partner commits; other agrees & \_\_\_\_ \\
A's information + B's skepticism & Info helps A; threatens B & \_\_\_\_ \\
Shared vs. separate finances & Different preferences on money & \_\_\_\_ \\
Career vs. caregiving & One works; one stays home & \_\_\_\_ \\
\midrule
\end{tabular}
\end{center}

\subsubsection*{1.3 Sequencing Strategy for Household}

\textbf{If Phase Mismatch (|φ\_A - φ\_B| > 0.3):}

\begin{quote}
Bring lower-phase partner to higher-phase partner's level \emph{before} joint decision.

\textbf{Step 1:} Identify lower-phase partner (e.g., B in PreAware, A in Triggered)

\textbf{Step 2:} Intervene with lower-phase partner separately using $I_{\text{AWARE}}$+$I_{\text{AWARE},\kappa}$ (information + feedback)

\textbf{Step 3:} Once both in similar phase (both Triggered), joint commitment ($I_{\text{HOW}}$) works

\textbf{Timeline:} 2-4 weeks typically sufficient to move one phase level
\end{quote}

---

% =============================================================================
% LEVEL 2: ORGANIZATION
% =============================================================================

\subsection*{Level 2: Organization --- Diagnostic Template}

\textbf{Unit:} Company/department/team with hierarchical structure

\textbf{Characteristic:} Nested portfolios; manager coordination; cascading messages

\subsubsection*{2.1 Organizational 9C Diagnostic}

\begin{center}
\small
\begin{tabular}{@{}lllll@{}}
\toprule
\textbf{Level} & \textbf{Actor} & \textbf{φ (Phase)} & \textbf{A (Aware)} & \textbf{Conflict?} \\
\midrule
Executive & CEO/CFO & \_\_\_\_ & \_\_\_\% & \\
Management & Dept heads & \_\_\_\_ & \_\_\_\% & A vs M? \\
Frontline & Team leads/staff & \_\_\_\_ & \_\_\_\% & M vs F? \\
\midrule
\end{tabular}
\end{center}

\subsubsection*{2.2 Org Portfolio Coherence Check}

\begin{quote}
Three-level portfolio design:
\end{quote}

\begin{center}
\begin{tabular}{@{}lll@{}}
\toprule
\textbf{Level} & \textbf{Portfolio} & \textbf{Check:Consistent?} \\
\midrule
Org (L2) & [Org-level initiatives] & YES / NO \\
Department & [Dept implementation] & Aligns with Org? \\
Team & [Daily execution] & Aligns with Dept? \\
\midrule
\end{tabular}
\end{center}

\textbf{If NO:} Identify misalignment and redesign one of the portfolios.

\subsubsection*{2.3 Manager Translation Role}

\begin{quote}
Manager must translate org-level message to team context:
\end{quote}

\begin{center}
\begin{tabular}{@{}lll@{}}
\toprule
\textbf{Org Says} & \textbf{Manager Should Translate To} & \textbf{Team Hears} \\
\midrule
``Empower employees'' & ``Here's your decision authority \& budget'' & Autonomy \\
``Reduce costs 20\%'' & ``We cut waste, keep quality'' & Shared sacrifice \\
``Digital transformation'' & ``New tools; training provided'' & Opportunity \\
\midrule
\end{tabular}
\end{center}

---

% =============================================================================
% LEVEL 3: NATIONAL
% =============================================================================

\subsection*{Level 3: National --- Diagnostic Template}

\textbf{Unit:} Regional/national government implementing multi-domain policies

\textbf{Characteristic:} Geographic heterogeneity in Ψ; multiple domains; political constraints

\subsubsection*{3.1 Regional Heterogeneity Assessment}

\begin{center}
\small
\begin{tabular}{@{}llccc@{}}
\toprule
\textbf{Context Dimension (Ψ)} & \textbf{Urban Core} & \textbf{Suburbs} & \textbf{Rural} & \textbf{Gap} \\
\midrule
Ψ\_K (Information quality) & 0.8 & 0.5 & 0.2 & 0.6 \\
Ψ\_F (Factor flexibility) & 0.9 & 0.6 & 0.3 & 0.6 \\
Ψ\_E (Economic development) & 0.8 & 0.6 & 0.2 & 0.6 \\
Ψ\_I (Institutional quality) & 0.7 & 0.6 & 0.3 & 0.4 \\
\midrule
\textbf{Average Ψ Gap} & & & & 0.55 \\
\midrule
\end{tabular}
\end{center}

\textbf{Interpretation:} Gap > 0.4 indicates need for region-specific portfolio variant

\subsubsection*{3.2 Multi-Domain Spillover Matrix (from Appendix MMM)}

\begin{quote}
Use Appendix MMM Section 2 to compute spillover matrix across domains (Health, Pension, Environment, Labor)

Target: Coherence Score > 0.25
\end{quote}

\subsubsection*{3.3 Decentralization Strategy}

\begin{quote}
\textbf{Central government sets:}
\begin{itemize}[nosep]
\item Goals (e.g., ``Reduce diabetes 20\%'')
\item Minimum budget allocation
\item Equity standards (underserved regions get support)
\end{itemize}

\textbf{Regions/cities design:}
\begin{itemize}[nosep]
\item Specific portfolio (emergent 10C intervention vectors)
\item Sequencing and timeline
\item Communication approach
\end{itemize}

\textbf{Result:} Coherent national strategy + local flexibility
\end{quote}

---

% =============================================================================
% LEVEL 4: GLOBAL
% =============================================================================

\subsection*{Level 4: Global --- Diagnostic Template}

\textbf{Unit:} Multi-nation coordination without global government

\textbf{Characteristic:} Massive context variation; interest conflicts; no enforcement

\subsubsection*{4.1 Global Context Assessment}

\begin{quote}
For each nation, estimate context dimensions:
\end{quote}

\begin{center}
\small
\begin{tabular}{@{}llccc@{}}
\toprule
\textbf{Dimension} & \textbf{Wealthy Nations} & \textbf{Middle-Income} & \textbf{Low-Income} & \textbf{Range} \\
\midrule
Ψ\_E (Economic dev.) & 0.9 & 0.5 & 0.2 & 0.7 \\
Ψ\_I (Institutions) & 0.8 & 0.4 & 0.2 & 0.6 \\
Ψ\_K (Information) & 0.8 & 0.5 & 0.2 & 0.6 \\
Ψ\_F (Factor flexibility) & 0.7 & 0.3 & 0.1 & 0.6 \\
\midrule
\end{tabular}
\end{center}

\textbf{Implication:} Same global portfolio will have radically different effectiveness in different countries.

\subsubsection*{4.2 Interest Compatibility Matrix}

\begin{quote}
Estimate γ (alignment) between nations' interests:
\end{quote}

\begin{center}
\begin{tabular}{@{}lccc@{}}
\toprule
\textbf{Nation Pair} & \textbf{Interest Alignment} & \textbf{γ Estimate} & \textbf{Conflict Risk} \\
\midrule
Wealthy vs. Wealthy & Often aligned & +0.5 & Low \\
Wealthy vs. Middle-Income & Mixed & 0.0–+0.2 & Medium \\
Wealthy vs. Low-Income & Often misaligned & -0.2–0.0 & High \\
Middle-Income vs. Middle-Income & Variable & -0.1–+0.3 & Medium \\
\midrule
\end{tabular}
\end{center}

\subsubsection*{4.3 Global Agreement Design}

\begin{quote}
Best practice for multi-nation coordination:

\textbf{Top-level goals:} Aspirational, global agreement (e.g., ``Net-zero by 2050'')

\textbf{Nationally determined contributions:} Each country designs own pathway based on context

\textbf{Equity mechanisms:} Transfer payments / technology support from wealthy to poor nations

\textbf{Accountability:} Transparent reporting; peer pressure; not enforcement

\textbf{Example:} Paris Climate Agreement structure
\end{quote}

---

% =============================================================================
% SECTION: CROSS-LEVEL SPILLOVER MATRIX
% =============================================================================

\subsection*{Cross-Level Spillover Effects}

\textbf{How do changes at one level affect outcomes at another?}

\begin{center}
\small
\begin{tabular}{@{}l|ccccc@{}}
\toprule
\textbf{FROM→TO} & L0 & L1 & L2 & L3 & L4 \\
\midrule
\textbf{L0 (Individual)} & — & ++ & + & ± & - \\
\textbf{L1 (Household)} & ++ & — & + & ± & - \\
\textbf{L2 (Org)} & + & + & — & ++ & + \\
\textbf{L3 (National)} & ± & ± & ++ & — & ++ \\
\textbf{L4 (Global)} & - & - & + & ++ & — \\
\midrule
\end{tabular}
\end{center}

\textbf{Legend:}
\begin{itemize}[nosep]
\item ++ = Strong spillover (plan for it)
\item + = Moderate spillover (consider it)
\item ± = Mixed or uncertain
\item - = Weak/negative spillover (minimal effect)
\end{itemize}

\textbf{Read as:} Row = Source level; Column = Target level

\textbf{Example:} L3 (National policy) → L1 (Household): ± (mixed effect — policy can help or hurt households depending on implementation)

---

% =============================================================================
% SECTION: PARAMETER COMPARISON TABLE
% =============================================================================

\subsection*{Parameter Comparison: How 9C Dimensions Vary by Level}

\begin{center}
\small
\begin{tabular}{@{}p{1.5cm}|ccccc@{}}
\toprule
\textbf{Parameter} & \textbf{L0} & \textbf{L1} & \textbf{L2} & \textbf{L3} & \textbf{L4} \\
& Individual & Household & Org & National & Global \\
\midrule
\textbf{Unit Size} & 1 & 2-5 & 50-5000 & 1M-100M & 8B \\
\textbf{Heterogeneity} & Low-Med & Med & High & Very High & Extreme \\
\textbf{Coordination} & Self & Joint & Hierarchical & Decentralized & Voluntary \\
\textbf{Timescale} & Weeks-Months & Months-Years & Quarters-Years & Years-Decades & Decades \\
\textbf{Enforcement} & Self & Social & Contractual & Legal & Weak \\
\textbf{Aggregation} & N/A & Couple-level & Org-wide & Population-level & Global \\
\midrule
\end{tabular}
\end{center}

---

% =============================================================================
% SECTION: IMPLEMENTATION CHECKLIST
% =============================================================================

\subsection*{Multi-Level Implementation Checklist}

\textbf{Before launching implementation at any level, verify:}

\begin{center}
\small
\begin{tabular}{@{}p{0.5cm}p{7cm}p{1.5cm}@{}}
\toprule
\# & \textbf{Checklist Item} & \textbf{Status} \\
\midrule
1 & Identified implementation level (L0-L4) & ☐ \\
2 & Completed 9C diagnostic for your level & ☐ \\
3 & Identified cross-level spillovers & ☐ \\
4 & Checked for phase/awareness mismatches & ☐ \\
5 & Estimated complementarities (γ matrix) & ☐ \\
6 & Designed portfolios at all affected levels & ☐ \\
7 & Verified coherence across levels (>0.25 if applicable) & ☐ \\
8 & Planned sequencing and timing by level & ☐ \\
9 & Identified manager/translator roles & ☐ \\
10 & Set up monitoring protocol for spillovers & ☐ \\
\midrule
\end{tabular}
\end{center}

---

% =============================================================================
% REFERENCES
% =============================================================================

\subsection*{References and Related Reading}

\textbf{Core Chapter:}
\begin{itemize}[nosep]
\item Chapter 23: Multi-Level Implementation (shows all levels with examples)
\end{itemize}

\textbf{Foundation (all levels):}
\begin{itemize}[nosep]
\item Chapters 16-22: Individual through national level details
\end{itemize}

\textbf{Cross-Level Coordination:}
\begin{itemize}[nosep]
\item Appendix MMM: Cross-domain spillovers (adapt for cross-level)
\item Appendix LLL: Domain templates (adapt for multi-level)
\end{itemize}


%%%%%%%%%%%%%%%%%%%%%%%%%%%%%%%%%%%%%%%%%%%%%%%%%%%%%%%%%%%%%%%%%%%%%%%%%%%%%%%

%%%%%%%%%%%%%%%%%%%%%%%%%%%%%%%%%%%%%%%%%%%%%%%%%%%%%%%%%%%%%%%%%%%%%%%%%%%%%%%
%%%%%%%%%%%%%%%%%%%%%%%%%%%%%%%%%%%%%%%%%%%%%%%%%%%%%%%%%%%%%%%%%%%%%%%%%%%%%%%
\section{The Fundamental Question}
\label{sec:nnn-fundamental}
%%%%%%%%%%%%%%%%%%%%%%%%%%%%%%%%%%%%%%%%%%%%%%%%%%%%%%%%%%%%%%%%%%%%%%%%%%%%%%%

\begin{tcolorbox}[
    colback=red!5!white,
    colframe=red!75!black,
    title={\textbf{Fundamental Question}}
]
\textbf{How does unknown integrate with and contribute to the
Evidence-Based Framework for understanding behavioral outcomes?}

This appendix addresses the core challenge of formalizing behavioral principles
within a rigorous theoretical framework while maintaining practical applicability.
\end{tcolorbox}


\section{Critical Foundations and Objections}
\label{sec:nnn-foundations}
%%%%%%%%%%%%%%%%%%%%%%%%%%%%%%%%%%%%%%%%%%%%%%%%%%%%%%%%%%%%%%%%%%%%%%%%%%%%%%%

\subsection{Objection 1: Scope Limitations}

\textbf{Objection:} The framework presented may have limited applicability
outside the specific domain contexts discussed.

\textbf{Response:} We acknowledge domain-specificity. The framework provides
general principles that should be calibrated to specific contexts. Cross-domain
validation remains an open research question.

\subsection{Objection 2: Measurement Challenges}

\textbf{Objection:} Key constructs may be difficult to measure reliably in
practice.

\textbf{Response:} We recommend using validated scales and instruments where
available, and developing domain-specific proxies where necessary. Sensitivity
analysis should assess robustness to measurement uncertainty.



%%%%%%%%%%%%%%%%%%%%%%%%%%%%%%%%%%%%%%%%%%%%%%%%%%%%%%%%%%%%%%%%%%%%%%%%%%%%%%%
%%%%%%%%%%%%%%%%%%%%%%%%%%%%%%%%%%%%%%%%%%%%%%%%%%%%%%%%%%%%%%%%%%%%%%%%%%%%%%%

%%%%%%%%%%%%%%%%%%%%%%%%%%%%%%%%%%%%%%%%%%%%%%%%%%%%%%%%%%%%%%%%%%%%%%%%%%%%%%%
\section{Key Results}
\label{sec:nnn-results}
%%%%%%%%%%%%%%%%%%%%%%%%%%%%%%%%%%%%%%%%%%%%%%%%%%%%%%%%%%%%%%%%%%%%%%%%%%%%%%%

The analysis yields the following key results:

\begin{enumerate}
    \item \textbf{Result 1:} [Primary finding with quantitative support]
    \item \textbf{Result 2:} [Secondary finding with implications]
    \item \textbf{Result 3:} [Practical application or boundary condition]
\end{enumerate}


\section{Summary}
\label{sec:nnn-summary}
%%%%%%%%%%%%%%%%%%%%%%%%%%%%%%%%%%%%%%%%%%%%%%%%%%%%%%%%%%%%%%%%%%%%%%%%%%%%%%%

This appendix has presented the key concepts and theoretical foundations.
The main contributions include:

\begin{enumerate}
    \item Formal definitions and theoretical framework
    \item Integration with the broader EBF structure
    \item Practical guidelines for application
\end{enumerate}

The content supports decision-making and behavioral analysis within the
Evidence-Based Framework, providing a foundation for further research
and practical implementation.


\section{Glossary of Symbols}
\label{sec:nnn-glossary}
%%%%%%%%%%%%%%%%%%%%%%%%%%%%%%%%%%%%%%%%%%%%%%%%%%%%%%%%%%%%%%%%%%%%%%%%%%%%%%%

For comprehensive definitions, see Appendix G (Master Glossary).

\begin{table}[htbp]
\centering
\caption{Key Symbols in Appendix NNN}
\begin{tabular}{lll}
\toprule
\textbf{Symbol} & \textbf{Meaning} & \textbf{Reference} \\
\midrule
$\gamma$ & Complementarity parameter & Appendix B \\
$\Psi$ & Context vector & Appendix V \\
$U$ & Utility function & Chapter 3 \\
\bottomrule
\end{tabular}
\end{table}


\section{References}
\label{sec:references}
%%%%%%%%%%%%%%%%%%%%%%%%%%%%%%%%%%%%%%%%%%%%%%%%%%%%%%%%%%%%%%%%%%%%%%%%%%%%%%%

See master bibliography (\texttt{bcm\_master.bib}) for complete citations.

\nocite{bcm_master}

% =============================================================================
% END OF APPENDIX NNN
% =============================================================================

\vspace{2cm}
\begin{center}
{\large \textbf{End of Appendix NNN: METHOD-MULTI-LEVEL}} \\
\vspace{0.5cm}
{\small \textit{Chapter 23 provides full guidance; use this appendix for quick reference at your implementation level}}
\end{center}
